\documentclass[11pt]{amsart}

\usepackage{amsmath,amssymb,amsthm}
\usepackage[margin=1in]{geometry}
\usepackage[hyphens]{url}
\usepackage{hyperref}
\hypersetup{colorlinks=true,linkcolor=blue,citecolor=blue,urlcolor=blue}

\theoremstyle{plain}
\newtheorem{theorem}{Theorem}
\newtheorem{proposition}[theorem]{Proposition}
\newtheorem{lemma}[theorem]{Lemma}
\newtheorem{corollary}[theorem]{Corollary}
\theoremstyle{definition}
\newtheorem{definition}[theorem]{Definition}
\theoremstyle{remark}
\newtheorem{remark}[theorem]{Remark}
\newtheorem{empresult}[theorem]{Empirical Result}

\newcommand{\R}{\mathbb{R}}
\newcommand{\C}{\mathbb{C}}
\newcommand{\Z}{\mathbb{Z}}
\newcommand{\N}{\mathbb{N}}

\title[Wirsching's Weak Covering Conjecture]{Wirsching's Weak Covering Conjecture: an extended
computation, a proven bound, and the structure of what resists it}

\author{Renato Augusto Tavares}

\begin{document}

\maketitle

\vspace{-2.5em}
\begin{center}
\normalsize Universidade Federal de Goi\'as \par
\texttt{rat@discente.ufg.br} \par
ORCID: \url{https://orcid.org/0009-0002-0196-3311}
\end{center}

\begin{abstract}
Wirsching's 1998 Weak Covering Conjecture asks for a sub-exponential $K(l)$ such that a large
enough set $R_{j-1,j} \subset \Z_3^*$, built from strictly decreasing exponent tuples, is forced to
cover every unit residue modulo $3^l$. We extend the exact computation of $j^*(l)$, the smallest
such covering budget, from $l=20$ to $l=23$, and compare four growth models for
$e(l) := j^*(l) - l\log_4 3$ against the extended table: a constant model is disfavored at every
statistic checked, but logarithmic, square-root, and slow-linear growth remain statistically
indistinguishable from one another. We give the best unconditional bound known to us,
$j^*(l) \le \tfrac{119}{104}l + O(1)$, via a mean-payoff-game relaxation of the covering budget. We
then show why the natural Fourier approach to a tighter bound cannot work: a uniform
square-root-cancellation argument is capped below even this bound, the bulk of the relevant
exponential sum's mass sits in an exponentially numerous population no sparse-exception argument
reaches, and the residues that actually fail to be covered sit at extremal, non-generic arithmetic
cells no averaging method sees. In their place we prove a genuine structural theorem: the set of
uncovered residues at one budget is contained in the union of its double and its quadruple at the
next, giving an explicit recursive upper bound on $j^*(l)$ that is empirically exact at every budget
checked. We falsify a natural strengthening of this theorem with an explicit witness-level
counterexample, prove three further exact laws governing the last residues to be covered, and name
precisely the two lemmas whose proof would either sharpen this bound further or show it cannot be.
The Weak Covering Conjecture itself remains open.
\end{abstract}

\noindent\textbf{Keywords:} Collatz map, 3-adic dynamics, covering conjecture, mean-payoff games,
exponential sums, Littlewood--Offord theory.

\noindent\textbf{2020 Mathematics Subject Classification:} 11B50, 11K70, 91A50, 11L07, 37A44.

\section{Introduction}
\label{sec:intro}

Wirsching's 1998 monograph~\cite{wirsching1998} studies the density of predecessor sets under the
$3n+1$ map through an averaging construction on the $3$-adic integers. A limiting Markov chain
built from that construction has an invariant density, and Wirsching's own route to positive
predecessor density passes through a covering question about a family of finite subsets of
$\Z_3^*$, the $3$-adic units. For integers $j,k\ge0$, write
\begin{equation}
\label{eq:Rjk}
R_{j,k} := \Bigl\{\, \sum_{i=0}^{j} 2^{a_i}3^i \ \Bigm|\ j+k \ge a_0 > a_1 > \cdots > a_j \ge 0
\,\Bigr\} \subset \Z_3^*.
\end{equation}
Each element of $R_{j,k}$ is built from a strictly decreasing sequence of $j+1$ ``digit positions''
chosen from $\{0,\dots,j+k\}$: the $i$-th chosen position, counting from the top, contributes
$2^{a_i}$ scaled by $3^i$. Wirsching's Weak Covering Conjecture asks for a sub-exponential $K(l)$
such that $|R_{j-1,j} \bmod 3^l| \ge K(l)\cdot 3^l$ already forces $R_{j-1,j}$ to cover every unit
residue modulo $3^l$. Equivalently, writing $j^*(l)$ for the smallest $j$ such that $R_{j-1,j}$
covers $(\Z/3^l\Z)^*$ exactly, the conjecture concerns the growth rate of
\begin{equation}
\label{eq:el}
e(l) := j^*(l) - l\log_4 3,
\end{equation}
the excess of the true covering budget over the leading-order rate a cardinality argument alone
predicts ($|R_{j-1,j}| = \binom{2j}{j} \sim 4^j$, so matching $2\cdot3^{l-1}$ residues needs
$j \sim l\log_4 3$ at leading order).

An earlier, unpublished manuscript of the author computed $j^*(l)$ exactly for $l=1,\dots,20$ by
direct enumeration and reported the resulting table without a growth-rate analysis. This paper does
three things. First, it extends that table to $l=23$ and gives a proper statistical comparison of
candidate growth models for $e(l)$ against the full range (Sections~\ref{sec:computation}
and~\ref{sec:growth}). Second, it gives the best unconditional upper bound on $j^*(l)$ we have
found, via a reformulation of the covering budget as a mean-payoff game
(Section~\ref{sec:mpg}). Third, and at greatest length, it reports a sustained attempt to close the
gap to the conjectured rate directly, first by the natural Fourier approach, and, once that approach
is shown to fail for a precise and checkable reason, by a direct study of the residues that actually
resist coverage (Sections~\ref{sec:gapA} and~\ref{sec:holdouts}). That second attempt does not
settle the conjecture. It produces a genuine theorem, a genuine falsification of a natural
strengthening of that theorem, and two precisely named open lemmas, each of which we believe is
more tractable than the conjecture itself and worth a dedicated attack in its own right.

\section{The computation}
\label{sec:computation}

Table~\ref{tab:jstar} gives $j^*(l)$ and $e(l)$ for $l=1,\dots,23$. The values for $l\le20$
reproduce the earlier manuscript's table exactly; they were independently recomputed here by a
from-scratch Rust reimplementation of the covering search (a bitset rotation dynamic program over
residues modulo $3^l$), validated against the original Python implementation at every $l\le20$
before being trusted for the extension. The values for $l=21,22,23$ are new.

\begin{table}[h]
\centering
\begin{tabular}{c|c|c||c|c|c||c|c|c}
$l$ & $j^*(l)$ & $e(l)$ & $l$ & $j^*(l)$ & $e(l)$ & $l$ & $j^*(l)$ & $e(l)$ \\ \hline
1 & 1 & 0.208 & 9 & 13 & 5.868 & 17 & 21 & 7.528 \\
2 & 4 & 2.415 & 10 & 15 & 7.075 & 18 & 22 & 7.735 \\
3 & 6 & 3.623 & 11 & 16 & 7.283 & 19 & 23 & 7.943 \\
4 & 7 & 3.830 & 12 & 17 & 7.490 & 20 & 24 & 8.150 \\
5 & 9 & 5.038 & 13 & 18 & 7.698 & 21 & 25 & 8.358 \\
6 & 10 & 5.245 & 14 & 19 & 7.905 & 22 & 26 & 8.565 \\
7 & 11 & 5.453 & 15 & 20 & 8.113 & 23 & 27 & 8.773 \\
8 & 12 & 5.660 & 16 & 20 & 7.320 & & & \\
\end{tabular}
\caption{$j^*(l)$ and $e(l) = j^*(l) - l\log_4 3$, $l=1,\dots,23$.}
\label{tab:jstar}
\end{table}

The single plateau in the table, $j^*(15)=j^*(16)=20$, is not an error: $l=15\to16$ is the only
increment among the $22$ steps from $l=1$ to $l=23$ at which the covering budget does not increase,
and it is exactly the event Section~\ref{sec:growth}'s plateau-frequency test measures.

$l=21$ fits in physical RAM once a memory-wasteful parallelization step in the original
implementation was corrected. $l=22$ required $500\,\mathrm{GiB}$ of swap, added to the machine
specifically for this computation, and took approximately $10{,}749\,\mathrm{s}$ ($\sim\!3$ hours)
with swap I/O as the dominant cost, against $748\,\mathrm{s}$ for $l=21$ without swap. $l=23$
required approximately $263\,\mathrm{GiB}$ of state and $375{,}615\,\mathrm{s}$ ($\sim\!104$
hours); the successful run had no interruption, so a checkpoint/resume mechanism added after an
earlier, unrelated policy-triggered kill of an interrupted attempt at the same level was never
exercised on it. Three attempts at $l=24$, made over the course of this project on the same
machine, each failed before completion (two to memory-policy kills, one to an apparent out-of-memory
condition with no surviving kernel log confirming it directly); $l=24$ is not attempted further, and
the table stops at $l=23$.

\section{The growth of $e(l)$}
\label{sec:growth}

We compare four candidate functional forms for $e(l)$'s growth against the tail $l=10,\dots,23$
($n=14$; the same range choice as the earlier manuscript, now with three more points): a constant
(stabilization), a logarithmic term, a square-root term, and a slow-linear term. Table~\ref{tab:aic}
gives the fit statistics.

\begin{table}[h]
\centering
\begin{tabular}{l|c|c|c|c}
model & $k$ & $\Delta\mathrm{AIC}$ & $\Delta\mathrm{BIC}$ & LOOCV RMSE \\ \hline
constant & 1 & 16.091 & 15.452 & 0.5200 \\
logarithmic & 2 & 1.512 & 1.512 & 0.2991 \\
square-root & 2 & 0.722 & 0.722 & 0.2909 \\
slow-linear & 2 & 0.000 & 0.000 & 0.2837
\end{tabular}
\caption{Model comparison on $e(l)$, $l=10,\dots,23$. $\Delta$AIC/BIC relative to the best-fitting
model (slow-linear).}
\label{tab:aic}
\end{table}

A constant model is disfavored sharply and consistently: $\Delta\mathrm{AIC}=16.09$,
$\Delta\mathrm{BIC}=15.45$, both growing at every $l$ added since $l=21$. The three growth models
remain statistically indistinguishable from one another by the conventional $\Delta\mathrm{AIC}>2$
rule of thumb, with slow-linear taking a narrow lead on both AIC and leave-one-out cross-validation;
the three candidate regressors ($\log l$, $\sqrt l$, $l$) are themselves highly correlated over this
range (pairwise correlation $\ge0.993$), so this ranking says more about how thin the margin
currently is than about which functional form is correct.

A direct combinatorial test tells the same story from a different angle. Among the $22$ increments
$j^*(l+1)-j^*(l)$ for $l=1,\dots,22$, exactly one is zero (the $l=15\to16$ plateau noted above). A
one-sided binomial test against the plateau probability a constant-$e(l)$ source would imply gives
$p=0.0426$, crossing the conventional $0.05$ threshold for the first time with the $l=23$ point
included; restricted to the same tail the model comparison above uses, the same test gives
$p\approx0.22$. We report both figures rather than only the more dramatic pooled one: $e(l)$ is an
exact, deterministic sequence, not a sample from a noisy process, so neither AIC/BIC nor this test
carries the sampling interpretation of a classical inference. Both are best read descriptively, as
``how well does a constant-$e(l)$ source fit the observed increment pattern,'' not as a formal
hypothesis test.

A constant $e(l)$ is not merely disfavored statistically; it is impossible.

\begin{proposition}[Unboundedness of $e(l)$]
\label{prop:lowerbound}
$e(l) \ge \tfrac14\log_2 l - O(1)$.
\end{proposition}

\begin{proof}
$R_{j-1,j}$ is, by the bijection between strictly decreasing exponent tuples in $\{0,\dots,2j-1\}$
and $j$-subsets of a $2j$-element set, of exact cardinality $\binom{2j}{j}$. Covering
$(\Z/3^l\Z)^*$, of size $2\cdot3^{l-1}$, requires the domain to be at least as large as the
codomain: $\binom{2j}{j} \ge 2\cdot3^{l-1}$ is necessary for $j \ge j^*(l)$. The standard asymptotic
$\binom{2j}{j} \sim 4^j/\sqrt{\pi j}$ inverts this inequality to
$j \ge l\log_4 3 + \tfrac14\log_2 l - O(1)$.
\end{proof}

This argument uses no arithmetic structure of $R_{j,k}$ beyond its cardinality, and a working
number theorist would write it down without pausing on it; we verified it three ways before
including it (by hand, against the exact table via direct integer search with zero violations at
every $l\le23$, and by an independent, cold rederivation from a second reviewer) precisely because
it is otherwise the kind of one-line claim that is easy to get an exponent wrong on. We record it as
a rigorous floor under the statistical picture above, not as progress toward the true rate.

\section{A mean-payoff-game upper bound}
\label{sec:mpg}

The covering budget admits an exact reformulation as a game of perfect information. A state is a
unit $z \in \Z_3^*$; a move of cost $d\ge0$ is legal at $z$ exactly when $2^d z \equiv 1 \pmod 3$,
and it replaces $z$ with $T_d(z) := (2^{d+1}z - 2)/3$. $j^*(l)$ is the value, over $l$ steps with
the modulus shrinking by one power of $3$ per step, of the max-over-$z$, min-over-policy total cost.

Restricting a policy to see only $z \bmod 3^k$ turns this infinite-information problem into a
finite, two-player \emph{mean-payoff game}: the minimizing player chooses a legal $d$ knowing only
$z \bmod 3^k$, an adversary then chooses the hidden next $3$-adic digit $\epsilon\in\{0,1,2\}$, and
the state updates to $T_d(z + 3^k\epsilon) \bmod 3^k$. By the classical theorem of Ehrenfeucht and
Mycielski~\cite{ehrenfeuchtmycielski1979}, every mean-payoff game on a finite graph has a value at
each state, securable by both players with memoryless (positional) strategies; here the value is
independent of the starting state and we write it $\rho_k$.

\begin{theorem}[Window relaxation bound]
\label{thm:mpg}
For every $k\ge1$ there is an explicit rational constant $C_k$ such that
$j^*(l) \le \rho_k\cdot l + C_k$ for all $l$, and $\rho_k$ is non-increasing in $k$.
\end{theorem}

$\rho_k$ and $C_k$ are computed exactly, as rationals, by a nested Howard strategy-improvement
solver, self-certified at every $k$ via a matching adversary lower bound extracted from the same
run. Table~\ref{tab:mpg} gives the values computed for $k=3,\dots,13$.

\begin{table}[h]
\centering
\begin{tabular}{c|c|c|c|c|c|c|c|c|c|c|c}
$k$ & 3 & 4 & 5 & 6 & 7 & 8 & 9 & 10 & 11 & 12 & 13 \\ \hline
$\rho_k$ & $2$ & $\tfrac53$ & $\tfrac32$ & $\tfrac32$ & $\tfrac75$ & $\tfrac{25}{19}$ & $\tfrac54$ &
$\tfrac{11}{9}$ & $\tfrac65$ & $\tfrac76$ & $\tfrac{119}{104}$
\end{tabular}
\caption{The window-relaxation value $\rho_k$, $k=3,\dots,13$.}
\label{tab:mpg}
\end{table}

At $k=13$, $\rho_{13} = 119/104 = 1.144231\ldots$, giving:

\begin{corollary}
\label{cor:mpgbound}
$j^*(l) \le \tfrac{119}{104}l + O(1)$.
\end{corollary}

This is, to our knowledge, the best unconditional bound on $j^*(l)$'s growth rate currently
established. The general technique, representing a covering-type combinatorial problem as a
mean-payoff game and computing a decreasing sequence of exact rational relaxation bounds, is not
new to this application: beyond Ehrenfeucht and Mycielski's foundational
theorem~\cite{ehrenfeuchtmycielski1979} itself, a close and recent methodological precedent proves
rationality and computability of the covering radius of a primitive sofic shift by essentially the
same construction, applied to a different object~\cite{meyerovitchyoung2026}. What is new here is
the specific application to $j^*(l)$, not the framing, and we cite both as precedent rather than
present the technique itself as a contribution of this paper.

\begin{remark}
The trailing behaviour of the computed $\rho_k$ sequence, $\rho_k \to L$ with $L\approx0.78$–$0.79$
under a naive $L+A/k$ extrapolation, is consistent with the conjectured true rate
$\log_4 3 = 0.7925\ldots$, but we emphasize that only the direction $j^*(l)\le\rho_k l + C_k$ is
proven here. The reverse inequality, $\inf_k \rho_k \le \limsup_l j^*(l)/l$, which would be needed
to conclude $\rho_k \to \log_4 3$ rather than merely $\rho_k \to$ some value $\ge \log_4 3$, is not
established, and we found evidence in the course of this work that the window relaxation may carry
an irreducible gap above the true rate. We flag this explicitly because an earlier internal draft of
this result asserted the reverse inequality without proof; it does not hold as stated.
\end{remark}

\section{The exponential sum, and why it resists a direct attack}
\label{sec:gapA}

Corollary~\ref{cor:mpgbound}'s bound, $1.1442\ldots\,l$, is well above the conjectured rate
$\log_4(3)\,l \approx 0.7925\,l$. The natural route to closing that gap is Fourier-analytic. By
orthogonality on $\Z/3^l\Z$, the number of tuples landing in residue class $z$ is
\begin{equation}
\label{eq:Nz}
N(z) = \frac1{3^l}\sum_{t=0}^{3^l-1} S(t)\,e\!\left(-\frac{tz}{3^l}\right), \qquad
S(t) := \sum_{j+k\ge a_0>\cdots>a_j\ge0} e\!\left(\frac{t}{3^l}\sum_{i=0}^j 2^{a_i}3^i\right),
\end{equation}
and a uniform bound $|S(t)| = o(|R_{j,k}|)$ for every nonzero $t$ would give $N(z)>0$ for every $z$
once $|R_{j,k}|$ is large enough, at a rate controlled directly by the exponent in the bound. This
approach does not work, for three separate and mutually reinforcing reasons.

\subsection{A hard ceiling on uniform cancellation}

\begin{proposition}
\label{prop:sqrtceiling}
Any proof of covering via a uniform square-root-type bound on $|S(t)|$, applied via the triangle
inequality across all $3^l-1$ nonzero frequencies, requires $j \ge 2\log_4(3)\,l \approx 1.585\,l$.
\end{proposition}

\begin{proof}
Write $T := |R_{j-1,j}| = \binom{2j}{j}$. Positivity of $N(z)$ for every $z$ via
$\sum_{t\ne0}|S(t)| < T$ combined with a uniform estimate $|S(t)| \lesssim \sqrt{\beta T}$ over the
$3^l-1$ nonzero frequencies forces $3^l\sqrt{\beta T} \lesssim T$, i.e.\ $T \gtrsim 3^{2l}$, i.e.\
$j \gtrsim 2\log_4(3)\,l$.
\end{proof}

Proposition~\ref{prop:sqrtceiling}'s threshold is \emph{worse} than Corollary~\ref{cor:mpgbound}'s
already-proven bound: any approach that only delivers generic square-root cancellation across every
frequency cannot even match what is already known by other means, let alone improve on it. This
rules out the naive version of the Fourier approach cleanly, and directs attention to whether a
finer, frequency-dependent treatment (a major/minor-arc split, in the language of the circle method)
can do better.

\subsection{The exceptional mass does not stay small}

Numerically, $|S(t)|$ is far from uniform: the frequencies of largest magnitude sit extremely close
to dyadic rationals of small denominator, $t/3^l \approx p/2^r$ for small $r$, with the quality of
approximation improving by almost exactly a factor of $3$ per unit increase in $l$, the signature of
an underlying connection to the classical effective irrationality measure of $\log_2 3$ via
transcendence theory (Baker's theorem on linear forms in logarithms). This connects the problem
directly to a 2011 discussion of Terence Tao's, applying Littlewood--Offord and Riesz-product
methods to an object built from the same $2^a3^b$-type exponential sums, in the closely related
setting of a hypothetical non-trivial Collatz cycle~\cite{tao2011blog}. Reindexing Tao's construction
identifies it as an exact boundary slice of $R_{j,k}$ itself, with the bottom exponent pinned at
$0$; his own conclusion, in his own words, is that ``once one uses the transcendence theory bound,
\dots{} Littlewood-Offord techniques are no longer of much use,'' and that he knows ``of no way to
get a nontrivial separation property between powers of $2$ and powers of $3$ other than via
transcendence theory methods.''

Whether the sparse major-arc resonances described above can nonetheless be excised and the
remainder controlled by an averaged bound is a separate, checkable question. It fails: at the
threshold $\eta \sim c/\sqrt T$ that actually governs the $L^1$ norm (rather than at a fixed
threshold, where exceptional frequencies do look sparse), the count of exceptional frequencies grows
genuinely exponentially in $l$, roughly threefold per unit increase, with the proportion exceeding a
fixed multiple of $\sqrt T$ settling to a stable positive constant. At $l=18$, summing every one of
the $8{,}014$ primitive coefficients exceeding even a generous fixed threshold of $0.001\cdot T$
contributes less than $132$ to an $L^1$ norm of $5226$: over $97\%$ of the mass lives in an
exponentially numerous population of near-average-magnitude coefficients, not in the visible spikes.
No sparse-exceptional-set repair of the naive approach is available.

\subsection{Coverage failure is an arithmetic, not a statistical, phenomenon}

The most direct test is to examine the residues that actually fail to be covered, one budget step
before covering succeeds. At every level checked, their count is tiny (single or low double digits)
against an ambient occupancy mean in the hundreds, sitting in cells whose local intensity, at every
resolution checked, would assign them probability below $e^{-20}$ under any locally homogeneous
model; a genuinely independent-points model at the same mean would already have covered the residue
set completely. A direct experiment confirms this is a phase phenomenon, not a magnitude phenomenon:
holding every $|S(t)|$ fixed while randomizing the phases of conjugate-symmetric pairs collapses the
observed extremal statistics to what a generic-phase model predicts; the actual, unscrambled phases
do not. Coverage failure at the relevant budget is decided by an exact, highly non-generic
arithmetic obstruction, invisible to any method that only controls magnitudes, averages, or low
moments of $S(t)$.

\begin{proposition}
\label{prop:1547}
$S(3^{l-1})/\binom{2m}{m} \to 1/\sqrt3$ as $l/m$ grows, for the family $R_{m-1,m}$.
\end{proposition}

\begin{proof}[Proof sketch]
$\sum_i 2^{a_i}3^i \equiv 2^{a_0} \pmod 3$, so $S(3^{l-1})/\binom{2m}{m}$ depends only on the
limiting distribution of the top exponent's parity, which converges to a fixed, computable mixture
of the two nontrivial cube roots of unity.
\end{proof}

Uniform decay across every nonzero frequency, the premise a naive statement of the Fourier approach
might suggest, is thus not merely difficult; it is false, exactly, at an explicit and easily
described family of frequencies.

\section{A theorem on the residues that resist coverage}
\label{sec:holdouts}

Write $H(l,j)$ for the set of units modulo $3^l$ \emph{not} covered by $R_{j-1,j}$, so that
$j\ge j^*(l)$ exactly when $H(l,j)=\emptyset$.

\begin{theorem}[Holdout doubling inclusion]
\label{thm:doubling}
For $j \ge l$, $H(l,j+1) \subseteq 2H(l,j) \cap 4H(l,j)$.
\end{theorem}

\begin{proof}
Shifting every exponent in an admissible tuple up by $1$ (respectively $2$) produces an admissible
tuple at budget $j+1$ whose value is exactly twice (respectively four times) the original. So if $z$
is representable at budget $j+1$ via a tuple all of whose exponents are $\ge1$ (respectively
$\ge2$), then $z/2$ (respectively $z/4$) is representable at budget $j$. Once $j\ge l$, prepending a
new smallest exponent equal to $j$ contributes $2^j3^j \equiv 0 \pmod{3^l}$ for $j\ge l$, so this
covers the remaining case and every $z \notin 2R_{j-1,j} \cap 4R_{j-1,j}$ failing to be represented
at budget $j+1$ must already fail at budget $j$ after halving or quartering.
\end{proof}

\begin{corollary}[Chain contraction]
\label{cor:chain}
If $x, 2x, 4x, \dots, 2^{t-1}x \in H(l,j+1)$ (all reduced modulo $3^l$), then
$x, 2x, \dots, 2^{t}x \in H(l,j)$.
\end{corollary}

\begin{corollary}[Run-length bootstrap]
\label{cor:bootstrap}
Writing $\mathrm{maxrun}(H)$ for the length of the longest doubling chain contained in $H$,
$j^*(l) \le j + \mathrm{maxrun}(H(l,j))$ for every $j\ge l$.
\end{corollary}

\begin{proof}
By Corollary~\ref{cor:chain}, $\mathrm{maxrun}(H(l,j))$ strictly decreases by at least $1$ at every
budget step until it reaches $0$, at which point $H$ is empty.
\end{proof}

\begin{empresult}
\label{er:bootstrap}
Corollary~\ref{cor:bootstrap} is empirically exact, $j^*(l) = j + \mathrm{maxrun}(H(l,j))$, at every
$(l,j)$ with $j\ge l+1$ we were able to compute, $l=5,\dots,20$ at every budget and $l$ up to $22$ at
the single budget $j=l+1$, including a case with $|H(l,j)| > 3\times10^6$.
\end{empresult}

The converse direction, that a doubling chain of the observed maximal length is not merely
sufficient but necessary for survival to that budget, is not proven, and is false as an unrestricted
statement: $1547 \in 2H(7,8)\cap4H(7,8)$, yet $1547 \notin H(7,9)$ (both of $1547$'s two
budget-$9$ witnesses span the full exponent range, the one case Theorem~\ref{thm:doubling}'s
argument does not directly control). We return to a restricted, still-open form of this converse in
Section~\ref{sec:further}.

A natural strengthening of Theorem~\ref{thm:doubling}, that every individual witness admits a
bounded-cost local repair, is also false.

\begin{proposition}[Failure of uniform local repair]
\label{prop:falsify}
There is no bounded $\Delta d$ such that every witness of a covered residue at Durfee depth $d$
lifts to a witness at depth at most $d+\Delta d$ for every residue newly required at the next
budget. Explicit witness-level exhaustive checks at two consecutive budget transitions find,
respectively, $100$ of $1458$ and $332$ of $4374$ children whose minimum achievable depth exceeds
their parent's minimum depth by more than one, up to three at the transitions checked; a
correspondingly exhaustive check of minimum-edit-distance witness repair finds distances up to $7$
at both transitions, against a distance of $1$ for the overwhelming majority.
\end{proposition}

This falsifies, in its originally proposed strong form, a natural conjecture this project had earlier
entertained (that coverage propagates via a uniform, bounded-cost local repair rule budget by
budget); its valid surviving content is Theorem~\ref{thm:doubling} and its corollaries above.

\section{Further structure of the last holdouts}
\label{sec:further}

The residues in $H(l,j^*(l)-1)$, the last to be covered, satisfy several further exact laws.

\begin{theorem}[Last-holdout parity]
\label{thm:parity}
Every element of $H(l,j^*(l)-1)$ is $\equiv1\pmod3$.
\end{theorem}

\begin{theorem}[Near-extinction bijection]
\label{thm:extinction}
$H(l,j^*(l)-1) = 2\cdot\{x \in H(l,j^*(l)-2) : x\equiv2\!\!\pmod3\}$, exactly.
\end{theorem}

Both theorems follow from an exact reparameterization of the covering problem by \emph{width}
$W := j+l-1$ rather than by budget $j$ and level $l$ separately: coverage at $(l,j)$, $j\ge l$,
depends on $(l,j)$ only through $W$, via a monotone family of subsets of $\{0,\dots,W\}$; at the
minimal width admitting a given residue, every admissible witness is forced to place its largest
exponent at the same fixed position, whose parity determines the residue's class mod $3$ directly.

\begin{corollary}[$\times4$-only inheritance]
\label{cor:times4}
A last holdout at level $l+1$ is, if it descends from a last holdout at level $l$ at all, related to
it by multiplication by $4$, never by $2$.
\end{corollary}

\begin{proof}
By Theorem~\ref{thm:doubling}, a last holdout at level $l+1$ pulls back via halving or quartering to
$H(l,j^*(l+1)-2)$. By Theorem~\ref{thm:parity} the level-$(l+1)$ holdout is $\equiv1\pmod3$, so its
half is $\equiv2\pmod3$ and its quarter is $\equiv1\pmod3$; only the quarter can coincide with
another level's own last-holdout class by Theorem~\ref{thm:parity} again.
\end{proof}

\begin{theorem}[Mod-$9$ two-class bound]
\label{thm:mod9}
Writing $J:=j^*(l)$, $H(l,J-1) \bmod 9 \subseteq \{4^J, 4^{J+1}\} \pmod9$.
\end{theorem}

Theorem~\ref{thm:mod9} follows from Theorem~\ref{thm:doubling}'s reduction modulo $9$ (where $2$ is
a primitive root, so residue classes correspond to elements of $\Z/6$ via discrete logarithm) applied
one step before extinction. Numerically, $H(l,J-1)$ occupies the single class $4^J$ at every level
we computed, $l=5,\dots,21$; the exclusion of the other logically possible class, $4^{J+1}$, is
equivalent to a single parity statement about the second-highest exponent of the extremal witness,
and is the sharpest concrete open question this part of the investigation produced.

\begin{proposition}[H-014 equivalence]
\label{prop:h014}
$\mathrm{maxrun}(H(l,l+1))$ bounded by a constant $C$, for all $l$, is logically equivalent to
$j^*(l) \le l + O(1)$, not merely sufficient for it.
\end{proposition}

\begin{proof}[Proof sketch]
Sufficiency is Corollary~\ref{cor:bootstrap} at $j=l+1$. For necessity, the same corollary read
downward gives $\mathrm{maxrun}(H(l,l+1)) \ge j^*(l)-l-1$ directly from the definition of maxrun as
a witness to non-covering.
\end{proof}

We flag this equivalence because it corrects an initial, mistaken framing of our own: the naive
expectation, that a random set of holdouts at the observed density would have doubling chains far
longer than the flat ceiling of $3$–$4$ we measured at every level $l=5,\dots,22$, is wrong. Once the
correct, rapidly declining density of $H(l,l+1)$ is used rather than its value at small $l$, the
expected longest chain under a matched-density independence model tracks the observed sequence
within about one unit at every level, including both the one rise and the one fall in the sequence.
$\mathrm{maxrun}(H(l,l+1))$ is not anomalously small; it is statistically typical for the set's
density. Proposition~\ref{prop:h014} converts this from a reassuring-looking numerical coincidence
into the correct reading: bounding $\mathrm{maxrun}(H(l,l+1))$ is exactly as hard as the rate bound
it would imply, not an easier route to it.

A genuine, well-defined reduction of the harder, converse direction remains. Writing $U(l,W)$ for
the covering family at width $W$ (Section~\ref{sec:further}'s reparameterization), the exact
identity $U(W+1) = U(W) \cup 2U(W) \cup \mathrm{Corner}(W+1)$ holds, where $\mathrm{Corner}(W+1)$
consists of the witnesses using both the top and bottom available exponent simultaneously. At every
width $W \ge 2l+1$ we were able to check exhaustively, $l=3,\dots,13$, $\mathrm{Corner}(W+1)$
contributes nothing new: every corner witness's value is already covered without it. If this
\emph{corner-redundancy} property holds at every sufficiently large width, in general, it would make
$j^*(l)$ computable from a single holdout set via Corollary~\ref{cor:bootstrap}'s converse, a real
practical gain for extending Table~\ref{tab:jstar} beyond $l=23$; it remains unproven.

\section{Discussion}
\label{sec:discussion}

Two named lemmas summarize what stands between this paper's results and a sharper theorem. Excluding
the class $4^{J+1}$ in Theorem~\ref{thm:mod9} would make the mod-$9$ single-class law, observed at
every level checked, a proven statement rather than an observation. Proving corner-redundancy at
every sufficiently large width, in the sense of Section~\ref{sec:further}, would make
Corollary~\ref{cor:bootstrap}'s bootstrap tight in general, not merely empirically exact at every
computed instance, and would resolve Proposition~\ref{prop:h014}'s question in the process. Neither
lemma is, on its face, as hard as the Weak Covering Conjecture itself; whether either is
substantially easier is, honestly, not yet known.

Section~\ref{sec:gapA}'s three findings are, we believe, of independent interest beyond this
specific covering problem. A finite exponential sum built from strictly decreasing, exponentially
weighted digit tuples is a natural object in several adjacent areas (finite-population $L$-statistics
at non-CLT frequencies and weights, Littlewood--Offord theory for structured coefficient sequences,
and the analytic side of the Collatz literature via Tao's own related discussion), and the specific
mechanism identified here, that generic cancellation caps out below what is already provable by
combinatorial means, that the exceptional mass does not concentrate on a sparse set at the scale
that matters, and that the actual obstruction is an extremal, non-generic arithmetic phenomenon
rather than a statistical one, may recur elsewhere in that literature under different names. We
would welcome a reader who recognizes it.

The Weak Covering Conjecture itself, and the exact asymptotic rate of $e(l)$, remain open.

\section{Code and data availability}
\label{sec:code}

The repository \url{https://github.com/faculdade/weak-covering-conjecture} is reserved for this
paper's code and data, organized by section, and will be populated before submission.

\subsection*{Acknowledgments}
AI assistance was used for textual review and translation.

\begin{thebibliography}{99}

\bibitem{wirsching1998}
G.~J. Wirsching, \emph{The Dynamical System Generated by the $3n+1$ Function}, Lecture Notes in
Mathematics, vol.~1681, Springer, Berlin, 1998.

\bibitem{ehrenfeuchtmycielski1979}
A.~Ehrenfeucht and J.~Mycielski, \emph{Positional strategies for mean payoff games}, Internat. J.
Game Theory \textbf{8} (1979), no.~2, 109--113.

\bibitem{meyerovitchyoung2026}
T.~Meyerovitch and A.~Young, \emph{Rationality and computability of the covering radius for sofic
shifts}, arXiv:2603.21449 (2026).

\bibitem{tao2011blog}
T.~Tao, \emph{The Collatz conjecture, Littlewood--Offord theory, and powers of $2$ and $3$},
blog post, terrytao.wordpress.com, August 25, 2011.

\end{thebibliography}

\end{document}
